\chapter{Ein schönes Exempel}
\quellseite{9}

An anderer Stelle der heutigen Nummer findet sich ein beachtenswerther Artikel über Mischehen zwischen Christen und Ungläubigen. Es werden aber heutzutage auch viele Ehen zwischen Rechtgläubigen und Falschgläubigen geschlossen, auch in unsern Kreisen, und auch darüber ist fortgesetzt Belehrung und Mahnung nöthig. Denn bei allen derartigen Mischehen ist die Gefahr vorhanden, daß der rechtgläubige Theil seine Kirche verläßt und mit dem falschgläubigen Gemahl in die falsche Kirche geht, oder daß beide, weil sie in kirchlicher Hinsicht nicht übereinstimmen, gleichgültig werden und überhaupt wenig mehr nach der Kirche fragen und ihre Kinder kirchlos aufwachsen oder in irgend eine Kirche gehen lassen. Wer mag die vielen Söhne und Töchter zählen, die unsere lutherische Kirche durch solche Mischehen mit Unirt-Evangelischen, Reformirten und Römisch-Katholischen verloren hat! Ein böses Exempel haben in dieser Hinsicht oft die Prinzessinnen an Fürstenhöfen gegeben, die, wenn sie sich verheiratheten, ohne Weiteres der Kirche ihres Mannes beitraten. Anders hielt es die fromme Luise von Preußen, die Gattin des Königs Friedrich Wilhelm III., der im Jahre 1817 in seinem Lande die verhängnißvolle Union zwischen der lutherischen und reformirten Kirche einführte. Die Königin Luise ging nicht mit ihrem Manne und ihren Kindern in der unirten Kirche zum heiligen Abendmahl, sondern blieb ihrer lutherischen Kirche treu. Als im Jahre 1876 ihr hundertjähriger Geburtstag in ganz Preußen gefeiert wurde, theilte der bekannte Kirchengeschichtschreiber Guericke darüber Folgendes mit: „Ein innigeres Familienleben als das am königlichen Hofe der Unvergeßlichen, ein treueres Band als das zwischen ihr und ihrem königlichen Gatten, eine gewissenhaftere Muttersorge als die ihre gegen die die reichsten Hoffnungen erweckenden Söhne, eine hingebendere Volks- und Vaterlandsliebe als die ihrige, ein demüthigeres und bescheideneres Zurücktreten vor sinnführenden Geistern in Kirche und Staat als das ihre ist schwer denkbar: und doch hat diese herrliche Fürstin, Gattin und Mutter selbst an dem Hofe der Hohenzollern und an dem Herzen eines Friedrich Wilhelm III. nicht umhin gekonnt, die heiligste und tiefste Gemeinschaft des innersten Lebens beziehungsweise daranzugeben, und ist, fern von den Altären des Hofes, des Gatten, der Kinder, still und geräuschlos in einem Gotteshause ihres väterlichen lutherischen Bekenntnisses (zu St. Nicolai in Berlin, unter der Seelsorge Propst Ribbecks) dem Tische des Herrn genaht.“ (Zeitschrift für die gesammte lutherische Theologie und Kirche, 38, 91.)
